% Appendix B: lotteries are outperformed by equal division (the Jensen argument).
% Referenced by S8; \input inside the appendices environment AFTER appendix-gap-rule
% when S8 is typeset. Uses the prop counter (renders as Proposition 2) and Lemma 1
% of Appendix A.
\section{Lotteries and equal division}\label{app:lottery}

Fix a round, a budget $B > 0$, and a pool $P \subseteq \{1, \ldots, n\}$ of $m$
researchers. The \emph{equal division} over $P$ grants each researcher in $P$ the
amount $B/m$. A \emph{lottery} over $P$ with $k \leq m$ winners selects $k$
researchers from $P$ uniformly at random and grants each winner $B/k$. Expected
output is evaluated as in Appendix~\ref{app:gap-rule}: given an allocation $g$, it is
$\sum_i \varphi_i(g_i)$, and for a random allocation, the expectation is taken over
the draw as well.

\begin{prop}[lotteries are outperformed by equal division]\label{prop:lottery}
For every pool $P$, budget $B$, and winner count $k \leq m$, the lottery's expected
output does not exceed the equal division's, with equality only if $k = m$. In
particular, a lottery over the whole population, at any winner count, produces at
most uniform funding's expected output.
\end{prop}

\begin{proof}
Researchers outside $P$ receive nothing under either scheme. For $i \in P$, the
lottery's grant to $i$ is a random variable $\tilde g_i$ equal to $B/k$ with
probability $k/m$ and $0$ otherwise, so $\mathbb{E}[\tilde g_i] = B/m$. By linearity,
the lottery's expected output is $\sum_{i \notin P} \varphi_i(0) + \sum_{i \in P}
\mathbb{E}[\varphi_i(\tilde g_i)]$. Since $\varphi_i$ is strictly concave
(Lemma~\ref{lem:return}(ii)), Jensen's inequality gives
$\mathbb{E}[\varphi_i(\tilde g_i)] \leq \varphi_i(\mathbb{E}[\tilde g_i]) =
\varphi_i(B/m)$ for every $i \in P$, with equality only if $\tilde g_i$ is
degenerate, that is, $k = m$. Summing over $i$ yields the claim; taking
$P = \{1, \ldots, n\}$ yields the comparison with uniform funding.
\end{proof}

\begin{remark}[what the proposition does and does not cover]
The proposition compares expected output within a fixed pool. It measures the effect
of the randomization alone: whatever pool review selects, replacing the lottery within
that pool by the equal division of the same money weakly increases expected output.
The proposition says nothing about the choice of pool, about the review ranking
within the pool that the equal division also ignores, or about anything a lottery
saves outside the model (review costs, proposal effort, bias). \Cref{sec:egalitarian}
measures the first two and cites the literature on the third.
\end{remark}
